% !TeX program = xelatex
%%%%%%%%%%%%%%%%%%%%%%%%%%%%%%%%%%%%%%%%%%%%%%%%%%%%%%%%%%%%%%%%%%%%%%%%%%%%%%%
%  扩列图 / kuolie -- 这是你唯一需要改样式的文件
%
%  编译:  latexmk main.tex        (或者  xelatex main.tex)
%  必须用 XeLaTeX。Overleaf 里请把编译器设为 XeLaTeX。
%%%%%%%%%%%%%%%%%%%%%%%%%%%%%%%%%%%%%%%%%%%%%%%%%%%%%%%%%%%%%%%%%%%%%%%%%%%%%%%
\documentclass{kuolie}

\kuoliesetup{
  %-- 强调色 ----------------------------------------------------------------
  % 内置: emerald / skyblue / red / pink / orange / nephritis / concrete
  %       darknight / sakura / lavender / cyan / mint
  accent      = skyblue,
  % 想用自定义颜色就改用下面这行(六位十六进制,不要写 #):
  % accenthex = EF4089,

  %-- 中文字体 --------------------------------------------------------------
  % auto      自动探测,优先 Noto Sans CJK SC,最后退回 TeX Live 自带的 Fandol
  % noto      Noto Sans CJK SC   (推荐,免费跨平台,Overleaf 自带)
  % sourcehan Source Han Sans SC (思源黑体,与 Noto 同源)
  % fandol    TeX Live 自带,零安装,但缺日文汉字(如「込」「樋」)
  % macos     PingFang SC        (仅 macOS)
  % windows   Microsoft YaHei    (仅 Windows)
  fontset     = auto,

  %-- 版式 ------------------------------------------------------------------
  paper       = a4,        % a4 / letter / square(4:5 社交平台方图) / poster
  cardstyle   = flat,      % flat 纯白极简 / soft 淡色底 + 白色圆角卡
  divider     = rule,      % rule 细线(awesome-cv 风) / dots / hearts / none
  columns     = 2,         % pairlist 的栏数
  highlight   = 2,         % 小标题前几个字用强调色(0 = 都不用)
  titlespread = 0.06,      % 小标题字距,0 = 不加

  % columnsep = 2.2em,     % 双栏间距
  % pairdelim = {\,{\color{kuolie-lighttext}\textperiodcentered}\,},
  % tagsep    = {\,{\color{kuolie-hairline}\textbar}\,},
}

%-- 头像(可选) --------------------------------------------------------------
% 选项: circle|rectangle, edge|noedge, left|right
\photo[circle, noedge, left]{assets/avatar-sample.pdf}

\begin{document}

%%%%%%%%%%%%%%%%%%%%%%%%%%%%%%%%%%%%%%%%%%%%%%%%%%%%%%%%%%%%%%%%%%%%%%%%%%%%%%%
%  身份卡 -- 顶部那一块
%
%  \profileline 的内容用【英文逗号】分隔,分隔符「｜」会自动加上。
%  中括号里是图标,可以留空 \profileline{...},图标名见 fontawesome5 手册。
%%%%%%%%%%%%%%%%%%%%%%%%%%%%%%%%%%%%%%%%%%%%%%%%%%%%%%%%%%%%%%%%%%%%%%%%%%%%%%%

\name{示例名字}
\alias{（可以叫我小示｜本质是一团占位符）}

\profileline[\faUser]{天秤座, 大学在读, INFP}
\profileline[\faHeart]{小裙子爱好者, 甜品无限激推, 废话文学, 高强度冲浪}
\profileline[\faInfoCircle]{偏向自娱自乐, 空间浓度较低, 婉拒同嫁 / CP 向}

% 可选参数: [C] 居中(默认) / [L] 左对齐 / [R] 右对齐
\makekuolieheader[L]


\kuoliesection{我推和坑}
%%%%%%%%%%%%%%%%%%%%%%%%%%%%%%%%%%%%%%%%%%%%%%%%%%%%%%%%%%%%%%%%%%%%%%%%%%%%%%%
%  我推和坑
%
%  \pair{作品}{角色} -- 按列填充,左栏先填满。条目之间不要留空行。
%  第二个参数留空 \pair{只写作品}{} 也可以。
%%%%%%%%%%%%%%%%%%%%%%%%%%%%%%%%%%%%%%%%%%%%%%%%%%%%%%%%%%%%%%%%%%%%%%%%%%%%%%%

\begin{pairlist}
  \pair{摇曳百合}{船见结衣}
  \pair{超时空辉夜姬}{主角三位（两位）*支持一体论}
  \pair{金牌得主}{岡崎 16 花}
  \pair{鸟}{红尾水鸲}
\end{pairlist}

\kuolienote{%
  一些其他最近喜欢的和准备看：%
  \kuolietags{摇曳百合（全般）, 金牌得主（全般）, 再见菈菈, 恋死,
              在生态箱里吃早餐, se 酱, 今日マチ子作品全般,
              艾莉欧与电气人偶, 比恋爱更青涩, 致彼时繁花}%
}


\kuoliesection{关于CP}
%%%%%%%%%%%%%%%%%%%%%%%%%%%%%%%%%%%%%%%%%%%%%%%%%%%%%%%%%%%%%%%%%%%%%%%%%%%%%%%
%  关于 CP -- 和「我推和坑」结构完全一样,换数据即可
%%%%%%%%%%%%%%%%%%%%%%%%%%%%%%%%%%%%%%%%%%%%%%%%%%%%%%%%%%%%%%%%%%%%%%%%%%%%%%%

\begin{pairlist}
  \pair{金牌得主}{依实依 / 光祈}
  \pair{摇曳百合}{结京}
  \pair{超时空辉夜姬}{辉彩八}
\end{pairlist}

\kuolienote{%
  实际上比较杂食，且左右不严格固定，好吃的我都会吃一口）%
  完全不介意对家，对家好吃我也吃！%
}


\kuoliesection{一些碎碎念}
%%%%%%%%%%%%%%%%%%%%%%%%%%%%%%%%%%%%%%%%%%%%%%%%%%%%%%%%%%%%%%%%%%%%%%%%%%%%%%%
%  一些碎碎念 -- 通栏,不分栏
%
%  \kuoliehl{...} 可以把一小段文字标成强调色。
%%%%%%%%%%%%%%%%%%%%%%%%%%%%%%%%%%%%%%%%%%%%%%%%%%%%%%%%%%%%%%%%%%%%%%%%%%%%%%%

\begin{kuolieitems}
  \item \kuoliehl{列外的占位 CP 相关真的麻烦您不要来}，拜托了 555
        （因为个人原因，没有说不好的意思!!）
  \item 婉拒大众雷点 / 把我当小粉丝 / 经常吵架 / 攀比内卷，先在这里致歉 ><
  \item 我很好说话!! 正好放假了，可以随时找我玩～小窗大欢迎 .\^{}\textunderscore\^{}.
\end{kuolieitems}

\kuolieclosing{好啦我说完了!! 感谢您的阅读～天天开心}


\makekuoliefooter{}{}{}

\end{document}
